% !TEX root = 00_instructivo.tex
\chapter{Transductores ópticos: Fototransistor CNY70 y sensor de color TCS3200}

\section*{Objetivos}
\begin{itemize}
    \item Implementar un portón óptico con CNY70 para detectar reflectancia en hojas de colores.
    \item Utilizar el TCS3200 para discriminar colores de forma programada.
\end{itemize}

\section*{Materiales y equipo}
\begin{itemize}
    \item Sensor CNY70 (reflectivo IR).
    \item Sensor TCS3200 (fototransistor de color).
    \item Arduino Uno.
    \item Mini protoboard de 9 canales.
    \item Hojas de colores (blanco, negro, rojo, verde, azul).
    \item Resistencias: 220 $\Omega$ (LED), 10 k$\Omega$ (pull-up opcional).
    \item Cables Dupont.
\end{itemize}

\section*{Instrucciones}
\begin{enumerate}
    \item Conecte el CNY70: LED IR a través de 220 $\Omega$ a 5V; fototransistor colector a 5V, emisor a A0 con resistor pull-down de 10k a GND.
    \item Conecte el TCS3200: VCC a 5V, GND a GND, S0/S1 a 5V (frecuencia 100\%), S2/S3 a pines digitales (para seleccionar color), OUT a pin digital (medición de frecuencia).
    \item Calibre el CNY70 midiendo voltaje analógico con hoja blanca y hoja negra.
    \item Mida la reflectancia para cada hoja de color (10 lecturas).
    \item Con el TCS3200, programe la lectura de componentes RGB y calcule el color dominante.
    \item Compare la capacidad de discriminación entre ambos sensores.
\end{enumerate}

\section*{Esquemático}
\begin{figure}[htbp]
\centering
\begin{circuitikz}[scale=0.8]
\draw (0,0) node[above]{Arduino Uno};
% CNY70
\draw (1,0.5) node[anchor=east]{5V} -- (2,0.5) node[anchor=west]{LED+};
\draw (1,1) node[anchor=east]{GND} -- (2,1) node[anchor=west]{LED-};
\draw (1,1.5) node[anchor=east]{A0} -- (2,1.5) node[anchor=west]{Emisor (pull-down 10k)};
% TCS3200
\draw (1,2.5) node[anchor=east]{5V} -- (3,2.5);
\draw (1,3) node[anchor=east]{GND} -- (3,3);
\draw (1,3.5) node[anchor=east]{D2} -- (3,3.5) node[anchor=west]{OUT};
\draw (1,4) node[anchor=east]{D3} -- (3,4) node[anchor=west]{S2};
\draw (1,4.5) node[anchor=east]{D4} -- (3,4.5) node[anchor=west]{S3};
% S0,S1 atados a 5V
\draw (3,2.5) -- (3.5,2.5) node[anchor=west]{S0,S1 a 5V};
\end{circuitikz}
\caption{Conexiones en mini protoboard. Use resistencias pull-down para el CNY70.}
\end{figure}

\section*{Preguntas de análisis}
\begin{enumerate}
    \item ¿Por qué el CNY70 no puede diferenciar el rojo del verde si ambos tienen reflectancia similar en IR?
    \item ¿Qué ventaja tiene el TCS3200 al usar filtros de color RGB?
    \item Proponga una aplicación industrial donde se necesite un sensor de bajo costo como el CNY70.
    \item ¿Cómo afecta la distancia sensor-superficie a la lectura del CNY70?
\end{enumerate}

\section*{Bibliografía}
\begin{itemize}
    \item CNY70 Datasheet. Vishay.
    \item TCS3200 Color Sensor Datasheet. TAOS.
    \item Arduino Color Sensing Tutorial – Adafruit.
\end{itemize}
